% Set from report/draft.md section 2. The prose is the maintainer's; this file is
% its transit into LaTeX and nothing in it was rewritten.
\section{Background}

\subsection{Stereo depth as inference}

Stereo matching admits a Bayesian reading. For each pixel, comparing a patch
against candidate positions in the other image produces a cost curve over
disparity --- a likelihood. A prior over depth or a smoothness assumption
supplies the rest, and the quantity of interest is the posterior: an estimate
with an uncertainty attached, not a point estimate.

Uncertainty is usually read from the cost curve. Curvature at the minimum gives
a local precision; the ratio between the best and second-best minima measures
ambiguity; agreement between left-to-right and right-to-left matching flags
pixels with no consistent correspondent. Estimates from multiple views are then
combined by precision weighting --- each measurement contributes in proportion
to its confidence.

Depth is not disparity. Recovering metric depth requires knowing the geometry,
and for a converged binocular system it requires knowing the fixation distance.
This is where the oculomotor framing enters the estimator rather than just the
camera model, and it is the origin of one of the report's more portable
findings.

One property of this pipeline must be stated, because it motivated the work and
because the record complicates it. A half-occluded pixel --- visible to one eye,
hidden from the other --- has no true correspondent. The matcher does not report
this; it finds the best available wrong match, and a sharp cost minimum at a
wrong disparity is genuinely sharp. The result is a confident, incorrect
estimate at exactly the geometry an active system might hope to resolve by
looking elsewhere.

The predecessor repositories measured this on photographs, across three
uncertainty readouts, and found variance in half-occluded regions \emph{lower}
than in matched ones --- confidence inverted where it matters most. On
random-dot stereograms the same matcher and the same statistic give the
opposite: variance roughly twice as high in half-occlusions, which is the safe
direction.

So the failure is not a stimulus-independent property of stereo confidence. It
appears on real imagery and not on synthetic dots, and whether that difference
is about realism or about contrast is unresolved --- a random-dot field is
contrast-uniform by construction and cannot separate them. That open question is
the sharper motivation: confidence fails where the scene is real, and the
synthetic fixture this project used sits on the safe side of the reversal.

\subsection{Concepts}

Four distinctions do most of the analytical work, and none was fully in place
when the project began.

\textbf{Allocation versus acquisition.} \emph{Where to spend limited processing
on data already sampled}, against \emph{where to point the sensor to obtain data
not yet sampled}. Most treatments of active vision conflate them. Separating
them is what makes the results here interpretable, and the separation has a
clean geometric form: a rotation about a fixed optical centre changes allocation
and cannot change acquisition.

\textbf{Unknowable is not uncertain.} A pixel with no valid correspondent is not
a pixel we are unsure about; it is one we cannot measure. The difference between
a binary \emph{can this be measured} and a continuous \emph{how well} turned out
to dominate.

\textbf{Foveal weighting versus foveal sampling.} Weighting says the fovea
contributes more to the posterior. Sampling says the periphery was never
resolved in the first place. On a uniformly sampled sensor the first can only
discard information already paid for. Only the second is what a retina does.

\textbf{Linearising about the fixation is a confidence mechanism.} Expanding the
depth relation about the fixated disparity is exact at the fovea and degrades
with the square of the offset. That degradation term does more than bound
approximation error: it inflates the variance of any estimate far from the
fixated depth, and thereby refuses precisely the confident-and-wrong matches
described above.

A fifth concept was not brought to the project but forced on it, and \S3.3 takes
it up: a stimulus is an instrument, and must be characterised before it is used
to measure.
